\documentclass[journal]{IEEEtran}

\usepackage{cite}
\usepackage{amsmath,amssymb,amsfonts}
\usepackage{algorithmic}
\usepackage{graphicx}
\usepackage{textcomp}
\usepackage{xcolor}
\usepackage{booktabs}
\usepackage{multirow}
\usepackage{array}
\usepackage{float}
\usepackage{hyperref}

\def\BibTeX{{\rm B\kern-.05em{\sc i\kern-.025em b}\kern-.08em
    T\kern-.1667em\lower.7ex\hbox{E}\kern-.125emX}}

\begin{document}

\title{A Formal Interface That Maps NARS Truth Values to Uncertainty-Conditioned Transformer Attention for Rare Disease Prediction}

\author{Arjun Vijay Prakash\thanks{Grade 10 Student, City Montessori School, Lucknow, Uttar Pradesh, India (e-mail: arjunv.prakash12345@gmail.com).}}

\maketitle

\begin{abstract}
Uncertainty-aware clinical prediction remains difficult when oncology data are heterogeneous, partially missing, and only moderately sized \cite{ma2019, ovadia2019}. We define a formal interface that maps transformer outputs to Non-Axiomatic Reasoning System (NARS) truth values and then uses NARS confidence to modulate attention during inference \cite{wang1995, wang2006, wang2013}. The method makes uncertainty an explicit, inspectable variable and specifies where failures can arise in variance estimation, case-level proposition extraction, and confidence gating. We instantiate the interface on TCGA-THCA to predict lymph node metastasis from a multi-modal fusion of merged clinicopathologic tables and a somatic-mutation MAF-derived binary gene matrix. After case-level merging and target filtering, 457 labeled cases remain and are split into 319 training, 69 validation, and 69 test cases. On the current held-out test split, a random-forest baseline reached an AUC of 0.6613, the baseline transformer reached 0.7252, and both the flat-confidence and NARS-gated variants reached 0.7261; the flat-confidence control achieved the best Brier score at 0.2118 and the best accuracy at 0.6812. The 95\% bootstrap AUC intervals overlap across all variants, so the benchmark should be read as a feasibility study of the uncertainty interface and its reliability semantics rather than a definitive performance gain.
\end{abstract}

\begin{IEEEkeywords}
neurosymbolic AI, NARS, transformer, thyroid carcinoma, lymph node metastasis, clinical prediction, uncertainty quantification, TCGA
\end{IEEEkeywords}

\section{Introduction}
\label{sec:intro}

Clinical prediction still encounters data-scarcity-like behavior even when a benchmark is larger than a boutique case-series dataset. Structured oncology tables contain heavy missingness, repeated diagnosis rows, and substantial post-processing requirements before they can support reliable prediction \cite{ma2019}. These issues increase estimation error and degrade calibration in neural predictors \cite{ovadia2019}. In this paper, we move from a hand-curated rare-disease cohort to TCGA-THCA, a larger and more standard thyroid oncology benchmark, and study lymph node metastasis as the binary endpoint.

Transformers provide flexible function approximators for clinical prediction tasks \cite{rasmy2021, azizi2021, li2019}. In low-data regimes, a transformer can output high predicted probabilities even when epistemic uncertainty is high \cite{ovadia2019, ghassemi2021}. Standard attention weights also lack an explicit variable that denotes evidential support for each attended feature \cite{goncalves2022}. This matters clinically because a black-box score without an explicit uncertainty channel is difficult to audit, defer, or hand off to a human decision-maker.

Non-Axiomatic Reasoning System (NARS) represents evidential support with a two-dimensional truth value $(f,c)$, where $f$ is frequency and $c$ is confidence \cite{wang1995, wang2006, wang2013}. NARS updates truth values by truth-value functions that depend on the amount and agreement of evidence.

We study the following research question: How can a transformer expose its uncertainty in a form that supports symbolic evidence aggregation, and how can aggregated evidential confidence control attention during inference?

This paper does not claim that NARS broadly outperforms neural baselines for thyroid oncology prediction. The claim is narrower: NARS truth values provide a principled, inspectable uncertainty variable that can be coupled to attention and evaluated transparently on a clinically relevant benchmark. The clinical motivation is therefore reliability and inspectability first, not ranking performance alone.

We contribute a formal interface specification. First, we define a mapping from a neural predictive distribution to a NARS truth value and a mapping from a truth value to an attention reweighting factor. Second, we specify the NAL truth-value operations required to aggregate multiple clinical propositions into a single conclusion under uncertainty. Third, we instantiate the interface for TCGA-THCA lymph node metastasis prediction as a worked application case.

\section{Related Work}
\label{sec:related}

\subsection{Neurosymbolic AI Integration}

Neurosymbolic AI seeks to combine neural networks' pattern recognition capabilities with symbolic systems' interpretable reasoning \cite{garcez2020, marra2021}. Recent surveys identify multiple integration strategies, including symbolic knowledge injection into neural networks, neural-symbolic hybrid architectures, and neurosymbolic reasoning systems \cite{bhuyan2024, colelough2025}.

Hilario \cite{hilario2013} provides an overview of neurosymbolic integration strategies and categorizes approaches by the interface between neural and symbolic components. The survey identifies tight coupling, loose coupling, and unified architectures as distinct paradigms with different trade-offs in computational efficiency and interpretability.

Colelough and Regli \cite{colelough2025} present a systematic review of neurosymbolic AI in 2024, highlighting advances in learning and inference. Their analysis identifies neurosymbolic integration for enhanced learning through fusion of symbolic reasoning with neural mechanisms, particularly for few-shot settings. The review emphasizes Logical Neural Networks that transform observations into logical facts as a promising direction.

Nawaz et al. \cite{nawaz2025} analyze the current state of neurosymbolic AI, emphasizing essential integration mechanisms for advanced cognitive systems. Their work identifies hybrid methodologies combining neural network adaptability with symbolic AI's logical rigor as critical for applications requiring both performance and explainability.

\subsection{Transformers for Clinical Reasoning}

Transformer architectures have transformed medical AI across multiple modalities. Rasmy et al. \cite{rasmy2021} developed Med-BERT, adapting BERT for structured electronic health records to predict disease outcomes. Their work demonstrates that domain-specific pretraining on clinical data significantly improves prediction accuracy over general-domain language models.

Li et al. \cite{li2019} conducted an empirical study on fine-tuning BERT-based models for large-scale EHR notes, establishing best practices for clinical representation learning. Their findings indicate that transformer models can capture clinical concepts from unstructured text but require careful handling of medical terminology.

Goncalves et al. \cite{goncalves2022} surveyed attention mechanisms for medical applications, analyzing whether current approaches address medical domain needs. Their review identifies that while attention mechanisms improve model interpretability, most approaches do not explicitly model clinical reasoning processes or uncertainty.

Dai et al. \cite{dai2025} proposed automated classification of clinical diagnoses in EHR using transformers with attention mechanisms and saliency mapping for clinical rationale. Their domain-adapted attention mechanism weights clinical concepts according to diagnostic relevance, improving both accuracy and interpretability.

\subsection{Few-Shot Learning for Rare Diseases}

Few-shot learning addresses the fundamental challenge of learning from limited examples, critical for rare disease diagnosis where training data is inherently scarce \cite{alsentzer2022}. Karjalainen et al. \cite{karjalainen2025} presented SHEPHERD, a few-shot geometric deep learning approach for rare disease diagnosis. SHEPHERD leverages biomedical knowledge graphs and graph neural networks to embed patients and diseases into a shared latent space, enabling diagnosis with zero or few labeled examples per disease.

Chen et al. \cite{chen2023} proposed dynamic feature splicing for few-shot rare disease diagnosis, enriching features of novel classes with few training samples. Their approach addresses the feature inadequacy problem in medical image analysis when rare disease examples are limited.

Quellec et al. \cite{quellec2020} developed few-shot learning for rare pathology detection in fundus photographs using unsupervised probabilistic models. Their method detects rare conditions given only a few examples, demonstrating the potential of few-shot approaches in ophthalmology.

Yoo et al. \cite{yoo2021} demonstrated few-shot learning for OCT diagnosis of rare retinal diseases, showing that deep learning models can learn from limited rare disease data when properly designed. Their feasibility study indicates that few-shot classification enables diagnosis of rare lesions with minimal training examples.

\subsection{Uncertainty Quantification in Clinical AI}

Uncertainty quantification (UQ) is essential for safe clinical deployment of machine learning models \cite{ghoshal2021}. Abad et al. \cite{abad2025} provide a comprehensive survey of UQ for machine learning in healthcare, analyzing applications across the ML pipeline from data preprocessing to model evaluation.

Bayesian Neural Networks (BNNs) learn distributions over network weights, providing principled uncertainty estimates \cite{blundell2015}. However, BNNs face computational challenges in training due to intractable posterior inference, often requiring approximations that affect calibration.

Gal and Ghahramani \cite{gal2016} demonstrated that dropout in neural networks approximates Bayesian inference, enabling uncertainty estimation with minimal computational overhead. This Monte Carlo dropout approach has been widely adopted for clinical uncertainty quantification.

Lakshminarayanan et al. \cite{lakshminarayanan2017} proposed deep ensembles for uncertainty estimation, training multiple models and aggregating predictions. Ensembles provide reliable uncertainty estimates but require training multiple models, increasing computational cost.

\subsection{NARS and Non-Axiomatic Logic}

Pei Wang's Non-Axiomatic Reasoning System provides a formal framework for reasoning under AIKR \cite{wang1995, wang2006}. NARS is based on non-axiomatic logic (NAL), a term logic with experience-grounded semantics where truth-values indicate evidential support rather than objective probabilities.

Wang and Awan \cite{wang2013} conducted a case study applying NAL in medical diagnosis, comparing the logic with binary logic and probability theory. Their analysis demonstrates NAL's ability to handle uncertain, inconsistent, and incomplete knowledge in diagnostic reasoning.

Wang \cite{wang2022} presents the theoretical foundations of NARS, defining intelligence as the capacity to adapt under insufficient knowledge and resources. The system operates with finite capacity in real-time, revising beliefs as new evidence arrives without assuming convergence to objective truth.

NARS truth-values are two-dimensional, consisting of frequency (positive evidence ratio) and confidence (evidential support amount) \cite{wang2013}. This representation enables explicit modeling of uncertainty and supports non-monotonic reasoning as new evidence arrives.

\subsection{Knowledge Graphs for Biomedical Reasoning}

Knowledge graphs integrate heterogeneous biomedical data for reasoning and prediction \cite{wu2023}. Abu-Salih et al. \cite{abusalih2023} surveyed healthcare knowledge graph construction, identifying state-of-the-art approaches and open challenges in the field.

Rajabi and Kafaie \cite{rajabi2022} reviewed knowledge graphs and explainable AI in healthcare, demonstrating how structured biomedical knowledge enables interpretable predictions. Their work highlights the importance of integrating domain knowledge for clinical decision support.

Jiang et al. \cite{jiang2024} proposed reasoning-enhanced healthcare predictions with knowledge graph community retrieval, combining KG community-level retrieval with large language model reasoning. Their approach leverages biomedical knowledge graphs to enhance prediction accuracy and explainability.

Chudasama et al. \cite{chudasama2025} developed interpretable hybrid AI by integrating knowledge graphs with symbolic reasoning in medicine. Their methodology combines complex biomedical data with ontologies and logical reasoning for context-aware AI.

\section{Theoretical Framework}
\label{sec:theory}

\subsection{NARS Truth-Value Semantics}

NARS employs a two-dimensional truth-value representation defined as a pair $(f, c)$, where $f$ represents frequency and $c$ represents confidence. For a statement $S$ with positive evidence $w^+$ and negative evidence $w^-$, the frequency is computed as:

\begin{equation}
f = \frac{w^+}{w^+ + w^-}
\label{eq:frequency}
\end{equation}

The confidence measures the amount of evidential support:

\begin{equation}
c = \frac{w^+ + w^-}{w^+ + w^- + k}
\label{eq:confidence}
\end{equation}

where $k$ is a system parameter representing the amount of evidence needed for maximum confidence \cite{wang2006, wang2013}. This formulation explicitly represents both the balance of evidence (frequency) and the amount of evidence (confidence).

In this paper, we assume the standard NARS default evidential horizon $k=1$. Writing total evidence as $W=w^+ + w^-$, Eq.~\eqref{eq:confidence} becomes
\begin{equation}
c = \frac{W}{W+1}.
\label{eq:confidence_k1}
\end{equation}
This explicit assumption ensures exact reproducibility of all confidence calculations used below.

NARS inference rules derive new truth-values from premises according to the available evidence. In \emph{strong deduction}, a premise $M \rightarrow P$ with truth-value $(f_1, c_1)$ and a premise $S \rightarrow M$ with truth-value $(f_2, c_2)$ yield $S \rightarrow P$ with truth-value:

\begin{equation}
f_{ded} = f_1 \cdot f_2
\label{eq:deduction_f}
\end{equation}

\begin{equation}
c_{ded} = \frac{c_1 \cdot c_2 \cdot f_1 \cdot f_2}{c_1 \cdot c_2 \cdot f_1 \cdot f_2 + k}
\label{eq:deduction_c}
\end{equation}

This product form ensures that deductive chains cannot increase frequency beyond the combined support of the premises; confidence is then attenuated by the system evidence horizon $k$ \cite{wang1995, wang2006, wang2013}.

For evidence aggregation (updating the same statement with multiple independent pieces of evidence), NARS uses \emph{revision}. Given two truth values $(f_1,c_1)$ and $(f_2,c_2)$ for the same statement, define
\begin{equation}
\omega_1 = \frac{c_1}{1-c_1}, \quad \omega_2 = \frac{c_2}{1-c_2}
\label{eq:revision_weights}
\end{equation}
then the revised truth value $(f_{rev},c_{rev})$ is
\begin{equation}
f_{rev} = \frac{\omega_1 f_1 + \omega_2 f_2}{\omega_1 + \omega_2}
\label{eq:revision_f}
\end{equation}
\begin{equation}
c_{rev} = \frac{\omega_1 + \omega_2}{\omega_1 + \omega_2 + 1}.
\label{eq:revision_c}
\end{equation}
This operation is required when combining a neural statement $S\rightarrow P\langle f_{neural},c_{neural}\rangle$ with a symbolic statement for the same link (or when multiple neural samples provide independent estimates) before the result is used to modulate attention.

\subsection{Transformer Attention Mechanisms}

Transformer models process input sequences through self-attention mechanisms that compute weighted representations of all positions. For an input sequence $X = [x_1, x_2, ..., x_n]$, the self-attention computes:

\begin{equation}
\text{Attention}(Q, K, V) = \text{softmax}\left(\frac{QK^T}{\sqrt{d_k}}\right)V
\label{eq:attention}
\end{equation}

where $Q$, $K$, and $V$ are query, key, and value matrices derived from $X$, and $d_k$ is the dimension of key vectors.

Multi-head attention extends this mechanism by applying multiple attention functions in parallel:

\begin{equation}
\text{MultiHead}(Q, K, V) = \text{Concat}(head_1, ..., head_h)W^O
\label{eq:multihead}
\end{equation}

where $head_i = \text{Attention}(QW_i^Q, KW_i^K, VW_i^V)$.

In clinical applications, attention weights can indicate which input features contribute most to predictions \cite{goncalves2022}. However, attention weights do not explicitly represent uncertainty or confidence in predictions, particularly problematic when training data is limited.

\subsection{Integration Architecture}

The proposed NARS-transformer hybrid integrates symbolic reasoning with neural pattern recognition through a bidirectional interface. The architecture consists of three main components:

\textbf{Neural Encoder}: The transformer encoder processes clinical input data (text, structured EHR, or imaging features) to produce contextualized representations $H = [h_1, h_2, ..., h_n]$.

\textbf{Symbolic Reasoner}: NARS maintains a knowledge base of clinical concepts and relationships represented as Narsese terms and statements. The reasoner performs inference using NAL rules, propagating truth-values through the knowledge graph.

\textbf{Integration Layer}: A bidirectional interface converts between neural representations and symbolic statements. Neural predictions are translated to NARS statements with appropriate truth-values, while NARS inference results guide neural attention.

The integration operates through the following mechanism. For a clinical prediction task, the transformer encoder first processes input data to generate representations. These representations are converted to NARS statements with truth-values based on prediction confidence:

\begin{equation}
(f, c) = \left(p, \frac{p(1-p) + \varepsilon}{p(1-p) + \varepsilon + \sigma^2}\right)
\label{eq:neural_to_nars}
\end{equation}

where $p$ is the predicted probability and $\sigma^2$ represents prediction variance from dropout or ensemble sampling. This mapping is a design choice that treats $p(1-p)$ as the intrinsic variance of a Bernoulli trial and $\sigma^2$ as epistemic variance, yielding $c$ as a reliability score that decreases when epistemic uncertainty is high. The small constant $\varepsilon>0$ prevents the degenerate edge case $c=0$ when $p\in\{0,1\}$ and provides a minimal ``evidence floor'' so that a perfectly certain predictor (very small $\sigma^2$) can still yield high confidence even at extreme probabilities.

Conceptually, Eq.~\eqref{eq:neural_to_nars} parallels a Beta-distribution method-of-moments construction (the conjugate-prior family for Bernoulli outcomes), where uncertainty controls an implied evidence mass; in our interface, however, $c$ remains an evidential-support quantity, not a probability.

NARS performs inference using these statements as evidence, deriving new conclusions through deduction, induction, and abduction. Because $c$ depends on the variance estimator, this interface can inherit known failure modes of uncertainty estimates under dataset shift (e.g., MC dropout or ensembles can become miscalibrated or underestimate epistemic uncertainty), which would inflate $c$ and over-trust weak evidence \cite{ovadia2019}.

A practical evaluation should therefore include a sensitivity analysis over $\varepsilon$ and the uncertainty estimator: (i) vary $\varepsilon$ over several orders of magnitude (e.g., $10^{-8}$ to $10^{-4}$) and measure how often clinical decisions change near decision thresholds; (ii) compare $\sigma^2$ from MC dropout \cite{gal2016} versus deep ensembles \cite{lakshminarayanan2017} and report discrimination, calibration (ECE/Brier), and decision-utility metrics under both in-distribution and shifted test conditions.

The resulting truth-values guide transformer attention in subsequent processing through the following normalized update:

\begin{equation}
\alpha_i^{\text{new}} = \frac{\alpha_i \cdot g(f_i, c_i)}{\sum_j \alpha_j \cdot g(f_j, c_j)}
\label{eq:attention_update}
\end{equation}

where $g(f, c) = c^{\gamma}$ with hyperparameter $\gamma > 0$. The function $g(f,c) = c^{\gamma}$ was chosen to use only confidence (not frequency) because $c$ encodes evidential reliability (how much support exists and how consistent it is), whereas $f$ primarily encodes the directional content of the proposition (supporting vs. negating). Using $c$ alone makes the attention gate behave like a reliability filter: propositions with uncertain support are downweighted regardless of whether they suggest a high or low probability outcome. Alternatives that include $f$ (e.g., $g(f,c)=c^{\gamma}f$) are possible but can over-suppress informative counter-evidence (low $f$) even when it is highly reliable (high $c$).

In practice, $\gamma$ should be selected on held-out validation data (e.g., grid search over $\gamma\in\{0.25,0.5,1,2,4\}$) to optimize a task-relevant objective such as validation NLL/Brier score and downstream decision utility; for clinical reporting, $\gamma$ can also be tuned to improve calibration of the final outputs after attention reweighting.

\subsection{Truth-Value Propagation for Data Scarcity}

The key advantage of NARS-transformer integration lies in truth-value propagation under data scarcity. When training samples are limited, neural predictions have high variance and poorly calibrated confidence. NARS provides explicit uncertainty representation through its two-dimensional truth-values.

For rare disease prediction with $n < 1000$ samples, the hybrid approach operates as follows. The transformer learns general patterns from available data, producing initial predictions. NARS incorporates these predictions as evidence with appropriately low confidence due to limited training data.

NARS inference rules propagate truth-values through clinical knowledge graphs, combining neural predictions with symbolic domain knowledge. For example, given evidence that "patient has symptom S" and "symptom S is associated with disease D," NARS derives "patient may have disease D" with computed truth-value reflecting both the neural prediction confidence and the strength of the symbolic association.

The confidence measure explicitly represents the amount of evidence supporting each conclusion \cite{wang2013}. When evidence is limited, confidence remains low, signaling uncertainty to clinicians. This explicit uncertainty representation is absent in pure neural approaches.

\section{Application to TCGA-THCA Multi-Modal Lymph Node Metastasis Prediction}
\label{sec:application}

\subsection{Benchmark Definition}

The current implementation uses the TCGA-THCA clinical tables together with somatic mutation calls to predict lymph node metastasis from structured multi-modal information. The raw repository dataset consists of five tab-separated clinical files, \texttt{clinical.tsv}, \texttt{exposure.tsv}, \texttt{family\_history.tsv}, \texttt{follow\_up.tsv}, and \texttt{pathology\_detail.tsv}, plus one or more somatic mutation MAF files in the same \texttt{data/} directory. The clinical files are not case-level as distributed, so the preprocessing pipeline first normalizes TCGA missing-value placeholders such as \texttt{--}, \texttt{'--}, \texttt{Not Reported}, and \texttt{Unknown}, filters the clinical table to the primary diagnosis rows, collapses repeated records to one row per \texttt{case\_submitter\_id}, and then left-joins the auxiliary tables onto that clinical base. After placeholder normalization, columns with more than 70\% missingness are removed before modeling.

For the genomic branch, MAF files are read with comment-aware parsing so \texttt{\#}-prefixed metadata headers are ignored safely. We derive \texttt{case\_submitter\_id} from the first 12 characters of \texttt{Tumor\_Sample\_Barcode}, filter \texttt{Variant\_Classification} to functionally relevant events (\texttt{Missense\_Mutation}, \texttt{Nonsense\_Mutation}, \texttt{Frame\_Shift\_Del}, \texttt{Frame\_Shift\_Ins}, \texttt{Splice\_Site}, \texttt{In\_Frame\_Del}, and \texttt{In\_Frame\_Ins}), and drop synonymous or clearly non-coding noise such as \texttt{Silent} and \texttt{Intron}. The filtered mutations are then reduced to the most frequently mutated genes and pivoted to a case-level binary mutation matrix, where repeated mutations in the same gene are collapsed with a max operator. This genomic matrix is left-joined onto the clinical base, and missing mutation indicators are filled with zero so clinically observed but genomically unprofiled cases remain usable.

The binary target is derived from \texttt{diagnoses.ajcc\_pathologic\_n}. We map $N0\rightarrow 0$ and $N1/N1a/N1b\rightarrow 1$, while cases with $NX$ or missing nodal status are excluded before splitting. On the current snapshot, this leaves 457 labeled TCGA-THCA cases from 507 merged case rows. The resulting benchmark is substantially larger and more standardized than the earlier repository cohort, while still preserving the missingness, modality-mismatch, and uncertainty challenges that motivate the NARS interface.

\subsection{Hybrid Model Architecture}

For TCGA-THCA lymph node metastasis prediction, we use a NARS-transformer hybrid with a case-level tabular input that concatenates clinical and genomic information. The selected feature set contains five numerical clinicopathologic variables (age at diagnosis, year of diagnosis, and three tumor-dimension measurements), thirteen categorical clinicopathologic variables including gender, race, ethnicity, pathologic T category, prior malignancy, morphology, laterality, tumor focality, residual disease, and extrathyroid extension, and a binary genomic mutation block derived from the filtered MAF files. These predictors are clinicopathologic retrospective benchmark variables, not a strictly non-invasive baseline-only panel. Numerical clinical variables are imputed with KNN imputation, categorical clinical variables with most-frequent imputation, and all imputers are fit on the training split only to avoid leakage. The genomic mutation indicators bypass the clinical standardization step and are passed directly as binary inputs. In the current repository state, the downloader-integrated pipeline resolves two local MAF files and expands the genomic branch to 50 binary mutation features, yielding a final 105-dimensional numeric input vector after preprocessing.

The transformer encoder processes these features through multiple attention layers, learning interactions among the encoded variables. The output layer produces an initial probability of lymph node metastasis.

The current repository version does not add a handcrafted THCA-specific symbolic rule base. Instead, it isolates the uncertainty-to-NARS and confidence-gated attention path: predictive means and variances are mapped into NARS truth values, per-feature attention variance is converted into confidence estimates, and those confidences reweight the attention distribution during inference.

\subsection{Inference Process}

Given a patient case, the hybrid model operates through the following inference process:

\textbf{Step 1}: The transformer encoder processes the encoded case features to generate contextualized representations and an initial lymph node metastasis probability $p_{neural}$.

\textbf{Step 2}: The neural prediction is converted to a NARS statement with truth-value $(f_{neural}, c_{neural})$ using Eq.~\eqref{eq:neural_to_nars}, where $f_{neural} = p_{neural}$ and $c_{neural}$ reflects prediction uncertainty from limited training data.

\textbf{Step 3}: Monte Carlo dropout is also used to summarize feature-level attention mean and variance. Those attention statistics are mapped into per-feature NARS confidence values that serve as reliability estimates for the attention update.

\textbf{Step 4}: The attention weights are reweighted by a tunable confidence gate of the form $c^\gamma$ and renormalized over features.

\textbf{Step 5}: The gated attention and token scores produce the final probability, while the resulting $(f, c)$ values provide explicit uncertainty estimates for inspection and export.

\subsection{Uncertainty Quantification}

The hybrid model provides two forms of uncertainty quantification. Aleatoric uncertainty arises from inherent randomness in the prediction task, while epistemic uncertainty stems from limited knowledge or data.

NARS confidence $c$ primarily captures epistemic uncertainty, decreasing as the amount of evidence supporting a conclusion decreases \cite{wang2013}. In the present TCGA benchmark, this confidence is driven by predictive and attention variance rather than by a manually curated disease rule base.

The frequency $f$ represents the balance of positive and negative evidence, providing a point estimate similar to neural network outputs. The combination $(f, c)$ enables nuanced uncertainty representation impossible with single probability values.

\section{Results}
\label{sec:results}

The central empirical question is whether NARS confidence improves discrimination or calibration on the TCGA-THCA multi-modal benchmark. On the 69-case held-out test split, the answer is again mixed rather than decisive: the flat-confidence and NARS-gated variants both reach an AUC of 0.7261, the baseline reaches 0.7252, the flat-confidence control achieves the best Brier score at 0.2118 and the best accuracy at 0.6812, and a random-forest baseline reaches 0.6613. Thus the transformer family still outperforms the tested classical tree baseline on this split, but the differences among transformer variants remain small. All bootstrap AUC intervals overlap, so the experiment supports feasibility and close comparability among the transformer variants rather than a confirmed superiority claim for any one of them.

To make the interface concrete, we also report an end-to-end numeric trace using the mappings from Eqs.~\eqref{eq:neural_to_nars} and \eqref{eq:attention_update}.
\par Consider a neural estimate and a symbolic rule that both support the same conclusion $S\rightarrow P$:
\begin{itemize}
\item Neural output: $p=0.8$, $\sigma^2=0.05$.
\item Symbolic rule: $(f_{sym},c_{sym})=(0.9,0.7)$.
\item Attention gain setting: $\gamma=2$.
\end{itemize}

From Eq.~\eqref{eq:neural_to_nars}, the neural statement has
\begin{equation}
(f_{neural},c_{neural})\approx(0.8,0.7619)
\end{equation}
(using $p(1-p)=0.16$ and taking $\varepsilon$ negligible relative to $p(1-p)$ and $\sigma^2$).

We then aggregate neural and symbolic evidence by revision (Eqs.~\eqref{eq:revision_f}--\eqref{eq:revision_c}). Using $\omega=c/(1-c)$ gives $\omega_{neural}\approx 3.2000$ and $\omega_{sym}\approx 2.3333$, so
\begin{equation}
(f_{rev},c_{rev})\approx(0.8422,0.8469).
\end{equation}

Finally, the confidence-gated attention gain is $g=c_{rev}^{\gamma}\approx 0.7173$, and the attention weights update by Eq.~\eqref{eq:attention_update}.

\begin{table}[H]
\caption{Worked simulation trace for revising neural and symbolic evidence and applying confidence-gated attention.}
\label{tab:simulation_trace}
\centering
\begin{tabular}{@{}lll@{}}
\toprule
Quantity & Value & How computed \\
\midrule
$(f_{neural},c_{neural})$ & $(0.8000,0.7619)$ & Eq.~\eqref{eq:neural_to_nars} with $(p,\sigma^2)=(0.8,0.05)$ \\
$(f_{sym},c_{sym})$ & $(0.9000,0.7000)$ & knowledge-base statement (given) \\
$(\omega_{neural},\omega_{sym})$ & $(3.2000,2.3333)$ & $\omega=c/(1-c)$ \\
$(f_{rev},c_{rev})$ & $(0.8422,0.8469)$ & Eqs.~\eqref{eq:revision_f}--\eqref{eq:revision_c} \\
Gain $g$ & $0.7173$ & $g=c_{rev}^{\gamma}$ with $\gamma=2$ \\
\bottomrule
\end{tabular}
\end{table}

We evaluated the implemented pipeline on the merged TCGA-THCA benchmark using the default experimental setting of 50 Monte Carlo dropout passes with $\gamma=2$, together with a transformer architecture of $d_{model}=64$, 4 heads, and 2 encoder layers. The labeled cohort contains 457 cases in total, split into 319 training cases, 69 validation cases, and 69 held-out test cases. In the current repository state, the multi-modal input contains 105 numeric dimensions after preprocessing: 5 scaled clinical numeric features, 50 binary genomic mutation indicators, and 50 one-hot encoded clinical categorical features. Table~\ref{tab:empirical_results} summarizes the held-out performance of that configuration together with bootstrap confidence intervals for AUC and expected calibration error (ECE) from a 10-bin equal-frequency reliability diagram. To address the standard tabular-learning objection that tree ensembles may outperform transformers, we also train a random-forest baseline on the exact same split and the exact same post-imputation encoded design matrix, selecting its hyperparameters by validation Brier score. To isolate whether the variance-to-confidence mapping itself matters, we additionally evaluate a flat-confidence control that uses the same gated architecture but assigns a constant feature confidence $c=0.5$ before attention renormalization. A gamma ablation is then run over $\gamma\in\{0.25,0.5,1,2,4\}$ using the same trained model and held-out split to examine how the confidence gate changes the final outputs.

\begin{table}[H]
\caption{Empirical test-set performance on the 69-case TCGA-THCA held-out split for lymph node metastasis prediction.}
\label{tab:empirical_results}
\centering
\begin{tabular}{@{}lcccc@{}}
\toprule
Variant & AUC [95\% CI] & Brier score & ECE & Accuracy \\
\midrule
Random forest & 0.6613 [0.5101, 0.7908] & 0.2200 & 0.1597 & 0.5942 \\
Baseline transformer & 0.7252 [0.5931, 0.8454] & 0.2118 & 0.1390 & 0.6667 \\
Flat-confidence transformer & 0.7261 [0.5934, 0.8468] & 0.2118 & 0.1390 & 0.6812 \\
NARS-gated transformer & 0.7261 [0.5957, 0.8460] & 0.2122 & 0.1403 & 0.6667 \\
\bottomrule
\end{tabular}
\end{table}

At the operating threshold of 0.5, the flat-confidence variant achieved the best accuracy at 0.6812 on the held-out set, compared with 0.6667 for the baseline transformer and NARS-gated variant and 0.5942 for the random forest. Relative to the baseline transformer, the flat-confidence and NARS-gated variants improve AUC by $0.0008$ each at the displayed precision, but the corresponding bootstrap intervals in Table~\ref{tab:empirical_results} overlap almost completely. On this test split, the ranking effect is therefore negligible and should not be interpreted as a statistically separated gain.

\begin{table}[H]
\caption{Gamma ablation on the 69-case held-out TCGA-THCA test split.}
\label{tab:gamma_ablation}
\centering
\begin{tabular}{@{}ccccc@{}}
\toprule
$\gamma$ & Baseline AUC & NARS-gated AUC & Baseline Brier & NARS-gated Brier \\
\midrule
0.25 & 0.7252 & 0.7261 & 0.21184 & 0.21187 \\
0.5 & 0.7252 & 0.7261 & 0.21184 & 0.21192 \\
1.0 & 0.7252 & 0.7261 & 0.21184 & 0.21200 \\
2.0 & 0.7252 & 0.7261 & 0.21184 & 0.21219 \\
4.0 & 0.7252 & 0.7294 & 0.21184 & 0.21264 \\
\bottomrule
\end{tabular}
\end{table}

The gamma ablation in Table~\ref{tab:gamma_ablation} shows that stronger confidence gating modestly improves the NARS-gated AUC and produces the best NARS-gated AUC at $\gamma=4$. In contrast, Brier score worsens gradually as $\gamma$ increases on this run, so the multi-modal benchmark now exhibits a clearer ranking-versus-calibration trade-off.

We also computed decision-curve analysis over threshold probabilities from 0.05 to 0.95 to assess clinical utility in terms of net benefit. On this benchmark, the transformer variants remain broadly similar, but the multi-modal run no longer produces the near-identical curves seen in the earlier clinical-only snapshot. Differences now appear at several intermediate thresholds, while the random-forest curve remains weaker over much of the clinically plausible range. The decision-curve result is therefore best read as showing close behavioral similarity within the transformer family together with modest threshold-level advantages for the gated variants at selected operating points.

Calibration was assessed with a reliability diagram using 10 equal-frequency bins and by expected calibration error. On this split, ECE was 0.1597 for the random forest, 0.1390 for the baseline model, 0.1390 for the flat-confidence control, and 0.1403 for the NARS-gated model. Thus, the flat-confidence control is marginally best on Brier score and is numerically best on ECE by a small margin over the baseline model in the saved metrics, while the NARS-gated model no longer improves the calibration summaries at the default $\gamma=2$ setting. The current empirical pattern is therefore mixed: the NARS gate preserves competitive discrimination, but the flat-confidence control is slightly stronger on the saved calibration metrics.

\section{Analysis}
\label{sec:analysis}

Eq.~\eqref{eq:neural_to_nars} makes confidence decrease as predictive variance increases, holding $p$ fixed. This behavior matches the goal of downweighting evidence that the model cannot support with stable predictions under repeated sampling.

Eq.~\eqref{eq:attention_update} enforces two invariants. It preserves $\sum_i \alpha_i^{\text{new}} = 1$. It sets $\alpha_i^{\text{new}} = 0$ when $c_i = 0$.

The empirical comparison in Table~\ref{tab:empirical_results} suggests that the main benefit of the gated variants in the current multi-modal run is not a large ranking gain. AUC is nearly unchanged across the three transformer variants, while the flat-confidence control improves thresholded accuracy and marginally improves the calibration summaries. In a clinical triage setting, this pattern is more naturally interpreted as a small shift in predictive reliability and decision-threshold behavior than as a stronger rank-ordering result. Because the bootstrap intervals still overlap heavily, none of these differences should be treated as a confirmed performance gain on the current test split.

The same experiment also shows a cautionary point. The Brier score differences are extremely small, with the baseline and flat-confidence models separated only in the fourth decimal place, so the calibration effect remains modest compared with the uncertainty around AUC. At the same time, the gamma ablation indicates that stronger gating can improve AUC while worsening Brier score, whereas the equal-frequency ECE estimate is slightly better for the baseline and flat-confidence models than for the NARS-gated model. This trade-off is consistent with the idea that confidence-based reweighting may help ranking without necessarily improving every calibration summary on a moderately sized held-out split.

The random-forest comparison also clarifies the benchmark narrative. On this split, the transformer family outperforms the tested tree baseline on AUC, Brier score, accuracy, and ECE, so the paper can defend the use of a transformer encoder empirically rather than only theoretically. This result still does not imply that the NARS gate itself is the main reason for the transformer advantage; it only shows that, under the current preprocessing and split, the implemented transformer pipeline is materially stronger than the selected random-forest comparator on the saved task summaries.

The decision-curve analysis sharpens the within-family interpretation. The gain from NARS guidance is not a uniform threshold-level utility improvement on this benchmark. Instead, the multi-modal run shows several modest threshold-dependent differences among the transformer variants, with the gated models tending to improve over the baseline at selected operating points rather than across the entire threshold range. Taken together, the ranking, calibration, and decision-utility results support feasibility and close behavioral similarity among the compared variants, but not a blanket claim of superiority across all threshold choices.

Compared with simply using MC dropout output as a soft confidence signal, the present interface still adds an explicit evidential variable in the shared $(f,c)$ space before attention is updated. In the current TCGA-THCA benchmark, the implementation isolates that variance-conditioned confidence path rather than adding a handcrafted disease-specific rule base. The flat-confidence control in Table~\ref{tab:empirical_results} sharpens this point: constant-confidence gating can slightly exceed the NARS-gated variant on AUC, but it does not preserve the same variance-conditioned evidential semantics and is slightly worse on Brier score. This differs from approaches such as SHEPHERD and other knowledge-graph-augmented predictors \cite{karjalainen2025, jiang2024}, which enrich representations or retrieval but do not expose a NARS-style evidential confidence variable that can be exported and fed back into attention. The contribution here is therefore not a larger knowledge graph or a different uncertainty estimator by itself, but a formal rule for turning uncertainty into an inspectable control signal in inference. In clinical terms, that is the paper's answer to the black-box problem: the model does not only emit a risk score, it also emits a confidence-like evidential state that can justify downweighting or deferring brittle predictions.

The interface can fail in three cases. First, the variance estimator can misrepresent epistemic uncertainty. MC dropout can under-estimate uncertainty under dataset shift \cite{ovadia2019}; for example, if a new cohort has a different pathology mix or staging distribution, $c$ may be inflated and the attention gate may over-weight uncertain predictions. Second, case-level proposition extraction can misalign source variables and encoded features; for example, if repeated diagnosis rows or missing-value placeholders are collapsed incorrectly, the wrong $c_i$ will be coupled to the attention update. Third, a future clinical rule base could contain conflicting or weak rules, producing diluted revised confidence that downweights informative features.

\section{Limitations}
\label{sec:limitations}

This paper reports an initial single-split empirical evaluation, but the evidence remains limited by dataset size and scope. The current run uses 457 labeled TCGA-THCA cases and a 69-case held-out test set, so the measured differences should still be treated as preliminary.

The 69-case test set still has limited statistical power to distinguish very small AUC gains at the 95\% level. With 34 positive and 35 negative test cases, the benchmark is more informative than the earlier small-cohort setup but still too small to support strong claims from changes on the order of a few thousandths of AUC.

The mapping in Eq.~\eqref{eq:neural_to_nars} depends on a variance estimate $\sigma^2$ from a chosen uncertainty quantification method \cite{gal2016, lakshminarayanan2017}. Different estimators change $c$ even when $p$ is fixed.

\subsection{Knowledge Engineering Cost}

The method can require a proposition inventory and a knowledge base with truth values when it is deployed with explicit symbolic rules. That requirement shifts effort from model training to knowledge engineering. The present TCGA-THCA benchmark intentionally does not add a handcrafted thyroid rule base, which keeps the experiment focused on the neural-to-NARS and confidence-gating interface. A fuller clinical deployment would still need expert curation of proposition wording and initial $(f_{sym},c_{sym})$ assignments.

This also creates an inter-rater problem. If two clinicians disagree on the appropriate truth value for a thyroid-cancer rule, the system has no automatic way to decide which prior is correct. In practice, such disagreements would need either consensus review, sensitivity analysis over rule weights, or explicit reporting of alternative rule sets.

The present experiment also used a single operating threshold and one stratified split. Additional resampling and external validation are still needed before stronger clinical claims can be made.

The current repository snapshot now includes multiple MAF files and supports the intended top-50-gene mutation panel through the downloader-integrated preprocessing path. Even so, the genomic coverage is still bounded by what has been downloaded locally, so future work should verify that the same results hold under larger cohort-level mutation exports and external validation cohorts.

\section{Discussion and Future Work}
\label{sec:planned_evaluation}

In addition to the empirical run reported in Sec.~\ref{sec:results}, the interface remains straightforward to validate more broadly in a reproducible way. The next stage of evaluation should test (i) the numerical behavior of the mappings, (ii) end-to-end predictive performance under stronger resampling and cohort variation, and (iii) clinician-facing interpretability.

\subsection{Synthetic Sanity Checks}
\label{sec:synthetic_checks}

Synthetic settings where $p$ and $\sigma^2$ are controlled remain useful for stress-testing the interface. For a grid such as $p\in\{0.01,0.1,0.5,0.9,0.99\}$ and $\sigma^2\in\{10^{-4},10^{-3},10^{-2},10^{-1}\}$, one can compute $(f,c)$ using Eq.~\eqref{eq:neural_to_nars} and the attention reweighting using Eq.~\eqref{eq:attention_update}.

Such experiments can report tables and plots showing how $c$ and $\alpha^{\text{new}}$ vary as a function of $(p,\sigma^2)$ for multiple settings of $\gamma$ (e.g., $\gamma\in\{0.5,1,2,4\}$), and quantify how often attention mass shifts across a fixed decision threshold as $\sigma^2$ increases.

\subsection{Dataset-Scarcity Ablation}
\label{sec:scarcity_ablation}

The present study already evaluates a transformer baseline versus a transformer+NARS interface on a TCGA-THCA lymph node metastasis setting using a structured clinical benchmark. A fuller follow-up evaluation should expand this to external thyroid oncology cohorts and should report sensitivity to feature selection, target definition, and missing-data handling.

To test label-scarcity behavior more rigorously, future work should subsample the training data to multiple sizes (e.g., $n\in\{50,100,250,500,1000\}$ when feasible), repeat each setting over multiple random seeds, and report:
\begin{itemize}
\item \textbf{Discrimination}: AUC and accuracy.
\item \textbf{Calibration}: Brier score and reliability diagrams.
\item \textbf{Decision utility}: net benefit using decision-curve analysis.
\item \textbf{Shift sensitivity}: an out-of-distribution split (temporal split, site split, or covariate-shifted split) to measure robustness under dataset shift.
\end{itemize}

\subsection{Ablation of $g(f,c)$ and $\gamma$}
\label{sec:g_ablation}

Ablating the attention-gating function and its hyperparameters would clarify how much of the observed AUC gain comes from the specific form of the gate. A useful comparison would include:
\begin{itemize}
\item $g(f,c)=c^{\gamma}$ (this paper),
\item $g(f,c)=c\cdot\left(1-2\lvert f-0.5\rvert\right)$ (confidence with a mild ``extremeness'' penalty), and
\item a variant where $c$ is derived from a conjugate-Beta interpretation of $(p,\sigma^2)$.
\end{itemize}

For each variant, $\gamma$ should be tuned on a held-out validation set and the same discrimination, calibration, and decision-utility metrics should be reported to show how design choices in the interface affect clinician-facing reliability.

\subsection{Error Analysis and Human-in-the-Loop Case Studies}
\label{sec:case_studies}

To evaluate interpretability, future work should select 10--20 representative cases (or simulated hard cases) and present the NARS-derived explanation trace: the main propositions used, their individual $(f,c)$ values, and how they combine into the final conclusion. This should also categorize common failure modes (variance underestimation, proposition extraction errors, and conflicting rules) and show how each failure mode manifests in the attention update.

Overall, the empirical result in this paper provides an initial proof of feasibility, while the experiments in Sec.~\ref{sec:planned_evaluation} remain the next systematic validation steps needed to determine whether the modest reliability and threshold-behavior differences observed here are robust across cohorts and clinical settings.

\section{Conclusion}
\label{sec:conclusion}

This paper defined a bidirectional interface between transformers and NARS that (i) converts neural predictive outputs into NARS truth values and (ii) uses NARS confidence to modulate transformer attention during inference. The key benefit is that evidential support becomes an explicit variable that can be aggregated symbolically and then reused as a control signal in neural inference. That explicit uncertainty channel is the main proposed mitigation for the clinical black-box problem: the model can expose what it does not confidently know instead of only returning a scalar risk score.

Concretely, the interface (a) maps $(p,\sigma^2)$ to $(f,c)$ so that confidence decreases as epistemic variance increases, and (b) maps $(f,c)$ to a normalized attention update via a tunable confidence gain. On the current TCGA-THCA multi-modal benchmark, the default parameter setting outperformed the tested random-forest baseline on all reported task summaries and produced closely matched results across the baseline, flat-confidence, and NARS-gated transformer variants: AUC values were 0.6613, 0.7252, 0.7261, and 0.7261 respectively, while the flat-confidence control achieved the best Brier score at 0.2118 and the best accuracy at 0.6812. Bootstrap AUC intervals overlapped across the transformer variants, so the result is best understood as preliminary evidence that the uncertainty interface is feasible on a larger standard oncology benchmark, rather than as a definitive clinical performance gain.

Future work should evaluate the interface under controlled data-scarcity and dataset-shift regimes, reporting discrimination, calibration, and decision-utility metrics, and comparing against stronger uncertainty-aware baselines such as gradient-boosted trees and ensemble-based tabular predictors that do not use symbolic aggregation. Future work should also repeat the analysis with broader TCGA-THCA mutation coverage so the intended larger genomic feature panel can be assessed directly, and should quantify knowledge-base construction effort and characterize failures due to mis-specified rules, conflicting evidence, and proposition extraction errors.

\section*{Acknowledgment}

The author thanks Pei Wang for developing the Non-Axiomatic Reasoning System (NARS). The author also acknowledges The Cancer Genome Atlas Research Network for making the TCGA-THCA clinical data publicly available through the NCI Genomic Data Commons.

\section*{Funding}

This work received no external funding.

\section*{Conflicts of Interest}

The author declares no conflicts of interest.

\section*{Ethics Statement}

This work is methodological and did not involve new studies with human participants or animals, nor the collection of patient data; therefore ethics approval and informed consent were not required.

\section*{Data and Code Availability}

The code for the implemented pipeline is available in the public repository \url{https://github.com/ArjunCodess/NGTA}. The benchmark analyzed in this work is derived from the TCGA-THCA clinical tables distributed through the NCI Genomic Data Commons; no new patient data were collected.

\section*{Author Biography}

\begin{IEEEbiography}{Arjun Vijay Prakash}
Arjun Vijay Prakash is a student at City Montessori School, Lucknow, Uttar Pradesh, India. His research interests include neurosymbolic AI, formal logic, and clinical informatics, with a focus on uncertainty-aware prediction for structured medical datasets.
\end{IEEEbiography}

\bibliographystyle{IEEEtran}
\bibliography{references}

\end{document}
